\chapter{Justificativa}

O desenvolvimento de um \textit{software} gratuito para dimensionamento de operações unitárias integra a teoria da engenharia química com ferramentas de programação e automação, democratizando o acesso a simulações para instituições de ensino, profissionais autônomos e pequenas empresas desprovidas de licenças comerciais. Isso favorece o aprendizado ativo, permitindo visualizações em tempo real de como parâmetros operacionais afetam o desempenho dos equipamentos. 

No contexto da Indústria 4.0, o domínio de linguagens de programação e bancos de dados é um diferencial para engenheiros químicos, dada a crescente digitalização de processos e a necessidade de análise de dados em tempo real.

\textit{Softwares} proprietários como \textit{Aspen Plus}\textregistered\ e \textit{HYSYS}\textregistered\ apresentam custo elevado e acesso restrito \cite{aspen2025, alves2022}, enquanto alternativas \textit{open-source} como o nacional \textit{DWSIM} \cite{dwsim2025} ainda apresentam limitações no acoplamento de módulos de forma didática e customizável. Planilhas em \textit{Excel}\textregistered\ ou rotinas em \textit{MATLAB}\textregistered, por sua vez, oferecem baixa modularização e carecem de interfaces amigáveis para simulação de processos complexos. Identifica-se, assim, uma lacuna em soluções abertas, escaláveis e focadas no aprendizado, que combinem rigor técnico com facilidade de uso. Este trabalho busca preencher essa lacuna com um sistema modular e extensível, servindo tanto ao ensino quanto à prática profissional.

Sendo de código aberto, espera-se que estudantes com formação interdisciplinar em engenharia química e desenvolvimento de \textit{software} possam contribuir para expandir e aprimorar o repertório da ferramenta, servindo como precursora para uso em maior escala.
